\documentclass[UTF8, scheme=chinese]{ctexbook}

\usepackage{xeCJK}
\xeCJKsetup{
  AutoFallBack = true,
  AutoFakeBold = true,
  PunctStyle = kaiming,
  CheckFullRight = true,
  CheckSingle = false,
  AllowBreakBetweenPuncts = true,
  CJKmath = true,
}

% 完整的中文章节格式设置
\ctexset{
    % 章设置
    % chapter = {
    %     format = \huge\bfseries\centering,
    %     nameformat = {},
    %     titleformat = {},
    %     name = {第,章},
    %     number = \chinese{chapter},
    %     aftername = \quad,
    %     beforeskip = 20pt,
    %     afterskip = 20pt,
    % },
    % 节设置
    section = {
        format = \Large\bfseries,
        nameformat = {},
        titleformat = {},
        name = {第,节},
        number = \chinese{section},
        aftername = \quad,
    },
    % 小节设置
    subsection = {
        format = \large\bfseries,
        nameformat = {},
        titleformat = {},
        name = {第,句},
        number = \chinese{subsection},
        aftername = \quad,
    },
    % 子小节设置
    subsubsection = {
        format = \normalsize\bfseries,
        nameformat = {},
        titleformat = {},
        name = {第,子小节},
        number = \chinese{subsubsection},
        aftername = \quad,
    }
}

\xeCJKDeclareSubCJKBlock{Ext-A}{ "3400 -> "4DB5 }
\xeCJKDeclareSubCJKBlock{Ext-B}{ "20000 -> "2A6DF }
\xeCJKDeclareSubCJKBlock{CircNum}{ "2460 -> "2473 }

\setCJKmainfont[%
BoldFont={FZYouHei_GBK 513B},   %
% Ext-A={仓耳今楷05-W04},
Ext-B={SimSun-ExtB},
CircNum={IosevkaTerm NFM}]
{NSimSun}
\newCJKfontfamily[li]\hmls{FZZJ-HMLSJF}
\newCJKfontfamily[kai]\fzxk{FZNewKai-Z03}
\newCJKfontfamily[song]\ssb{SimSun-ExtB}

\usepackage{xcolor}
\usepackage{pifont}

% 引入 tikz 及相关路径变形库（用于画波浪线）
\usepackage{tikz}
\usetikzlibrary{decorations.pathmorphing}

% 自定义双波浪线命令 \uuwave
\newcommand{\uuwave}[1]{%
  \tikz[baseline=(char.base)]{
    % 将文字放入一个无边距的节点中
    \node[inner sep=0pt, outer sep=0pt] (char) {#1};
    % 画第一条波浪线 (距离文字稍近)
    \draw[decorate, decoration={snake, amplitude=0.3mm, segment length=2mm}]
      ([yshift=-1.5pt]char.south west) -- ([yshift=-1.5pt]char.south east);
    % 画第二条波浪线 (在第一条下方)
    \draw[decorate, decoration={snake, amplitude=0.3mm, segment length=2mm}]
      ([yshift=-3.0pt]char.south west) -- ([yshift=-3.0pt]char.south east);
  }%
}

% 引入下划线宏包，normalem 参数防止破坏原有的 \emph{} 斜体命令
\usepackage[normalem]{ulem}

% 引入颜色宏包（可选，用于让中文标注颜色变浅一点）
\usepackage{xpinyin}

% 设置段落缩进和行距
\setlength{\parindent}{0pt}
\linespread{1.5}

% ==========================================
% 自定义语法标注命令 \markword
% 用法: \markword[线型命令]{英文内容}{中文标注}
% 默认线型为单实线 \uline。如果不想要线，可以传入 \mbox
% ==========================================
\newcommand{\markword}[3][\uline]{%
  \begin{tabular}[t]{@{}c@{}}%
    #1{#2} \\[-0.2ex]% 调整英文和中文之间的垂直间距
    \small\color{darkgray} #3% 中文标注字号设小，颜色设为深灰
  \end{tabular}%
}

\usepackage{mathabx}

\title{\fzxk {\huge 雅思语法学习笔记}}

\begin{document}
\maketitle

\tableofcontents

\chapter{词性}

\section{名词和冠词}

\subsection{复合名词}

\markword[\uuline]{\textbf{Britain}}{主语 1}
\markword{had produced}{谓语}
\markword[\uwave]{textiles}{宾语}
\markword[\mbox]{(like wool, linen and cotton),}{后置定语 (修饰 textiles)}
\markword[\mbox]{[for hundreds of years],}{时间状语}
\vspace{0.5cm} % 手动换行并增加行距
\markword[\mbox]{but}{连词}
\markword[\mbox]{[prior to the Industrial Revolution],}{时间状语}
\markword[\uuline]{the British textile business}{主语 2}
\markword{was}{系动词}
\vspace{0.5cm}
\markword[\uuwave]{a true ``cottage industry'',}{表语}
\markword[\mbox]{[with the work performed in small workshops or even homes by}{伴随状语}
\vspace{0.5cm}
individual spinners, weavers and dyers].

\subsection{可数名词变复数（规则变化）}

\markword[\mbox]{There's been}{there be 结构}
\markword[\uuline]{\textbf{some interesting research}}{主语}
\markword[\mbox]{{\large \textcolor{red}{\(\lgroup\)}}}{定语从句(修饰 research)开始}
\markword[\mbox]{which}{引导词/主语}

\vspace{0.5cm}

\markword{has found}{谓语}
\markword[\mbox]{\textcolor{blue}{\(\lgroup\)}}{宾语从句开始}
\markword[\mbox]{that}{引导词}
\markword[\uuline]{tardigrades}{主语 1}
\markword{are capable of}{谓语 1}
\markword[\uwave]{surviving radiation and very high pressure,}{宾语 1}

\vspace{0.5cm}

\markword[\mbox]{and}{连词}
\markword[\uuline]{they}{主语 2}
\markword{are also able to withstand}{谓语 2}
\markword[\uwave]{temperatures}{宾语 2}
\markword[\mbox]{(as cold as -200 degrees centigrade),}{后置定语，修饰 temperatures}

\vspace{0.5cm}

\markword[\mbox]{or}{连词}
\markword[\uwave]{highs}{宾语 3}
\markword[\mbox]{(of more than 148 degrees centigrade),}{后置定语，修饰 highs}
\markword[\mbox]{(which is incredibly hot).}{定语从句}
\markword[\mbox]{\textcolor{blue}{\(\rgroup\)}}{宾语从句结束}
\markword[\mbox]{{\large \textcolor{red}{\(\rgroup\)}}}{定语从句结束}

\subsection{可数名词变复数（不规则变化）}

\markword[\mbox]{[In the 1960s],}{时间状语}
\markword[\uuline]{the astronomer Gerald Hawkins}{主语}
\markword{suggested}{谓语}
\markword[\mbox]{\large\textcolor{blue}{\(\lgroup\)}}{宾从开始}
\markword[\mbox]{that}{引导词}

\vspace{0.5cm}

\markword[\uuline]{the cluster of megalithic stones}{主语}
\markword[\mbox]{operated}{谓语}
\markword{[as a form of calendar],}{方式状语}

\vspace{0.5cm}

\markword[\mbox]{[with different points corresponding to astronomical phenomenal]}{伴随状语}

\vspace{0.5cm}

\markword[\mbox]{<such as solstices, equinoxes and eclipses>}{同位语（修饰 phenomenal）}
\markword[\mbox]{(occurring at different times of the year).}{定语（修饰同位语）}
\markword[\mbox]{\large\textcolor{blue}{\(\rgroup\)}}{宾从结束}


\end{document}

% Local Variables:
% TeX-engine: xetex
% TeX-master: t
% End:

% LocalWords:  AutoFallBack AutoFakeBold PunctStyle kaiming CheckFullRight GBK
% LocalWords:  CheckSingle AllowBreakBetweenPuncts CJKmath BoldFont FZYouHei li
% LocalWords:  NSimSun IosevkaTerm NFM FZZJ HMLSJF FZNewKai kai SimSun ExtB
% LocalWords:  CircNum xetex LocalWords
